% La legge di Stevino
% OrbBits LaTeX conventions: see docs/style/latex.md
% Several languages: `orbbits build` defines \orbLocale per localised output; add
% \providecommand{\orbLocale}{it} here and switch babel and text on it (docs/localization.md).
\documentclass[11pt,a4paper]{article}
\usepackage[T1]{fontenc}
\usepackage[utf8]{inputenc}
\usepackage[italian]{babel}
\usepackage{lmodern}
\usepackage{amsmath,amssymb}
\usepackage[a4paper,hmargin=2cm,vmargin=1.7cm]{geometry}
\usepackage{graphicx}
\usepackage{xcolor}
\usepackage{enumitem}
\usepackage{microtype}
\usepackage[hidelinks]{hyperref}
\usepackage{tikz}
\usetikzlibrary{arrows.meta,decorations.pathreplacing}

\hypersetup{pdftitle={La legge di Stevino}}

% The handout is a tuned one-pager: block paragraphs, no page number.
\setlength{\parindent}{0pt}
\setlength{\parskip}{6pt}
\pagestyle{empty}

\begin{document}

{\LARGE\bfseries La legge di Stevino}\par
\vspace{2pt}
{\large Una colonna di fluido in equilibrio}
\vspace{8pt}

Consideriamo un \textbf{fluido in quiete}, di densità uniforme $d_f$.
Immaginiamo una colonna di altezza $h$ e base di area $S$: vogliamo
ricavare la pressione $p$ alla base, nota quella superiore $p_0$.
Assumiamo costante $g$. Nel disegno il fluido è un liquido esposto
all'aria, quindi $p_0=p_{\mathrm{atm}}$.

\vspace{5pt}
{\large\bfseries 1. Quali forze tengono ferma la colonna?}
\par\vspace{3pt}
\begin{minipage}[c]{0.36\textwidth}
\centering
\begin{tikzpicture}[x=1cm,y=1cm,>=Latex,font=\small]
  \fill[black!4] (-1.1,-0.45) rectangle (3.1,4.4);
  \fill[black!10] (0,0) rectangle (1.8,4.4);
  % Superficie libera: continua sopra la colonna, ondulata ai lati.
  \draw[thick] (-1.1,4.4) .. controls (-0.8,4.49) and (-0.5,4.31) .. (0,4.4)
    -- (1.8,4.4) .. controls (2.2,4.49) and (2.6,4.31) .. (3.1,4.4);
  \node[above=6pt,font=\footnotesize] at (0.9,4.4) {aria: $p_0=p_{\mathrm{atm}}$};
  \node[font=\footnotesize,rotate=90] at (-0.75,2.9) {liquido};
  \draw[dashed] (0,0) -- (0,4.4);
  \draw[dashed] (1.8,0) -- (1.8,4.4);
  \draw[very thick] (0,0) -- (1.8,0);
  \node[left=5pt] at (0,0) {$S$};
  \draw[->,very thick] (0.9,4.4) -- (0.9,3.7)
    node[midway,right=5pt] {$F_0$};
  \fill (0.9,2.2) circle (1.3pt);
  \draw[->,very thick] (0.9,2.2) -- (0.9,1.4)
    node[midway,right=5pt] {$F_P^{f}$};
  \draw[->,very thick] (0.9,0) -- (0.9,1.2)
    node[midway,left=5pt] {$R$};
  \draw[decorate,decoration={brace,amplitude=5pt}] (2.3,4.4) -- (2.3,0)
    node[midway,right=8pt] {$h$};
  \node[align=center,font=\footnotesize,text width=5.7cm] at (1,-0.95)
    {Le forze che agiscono sulla colonna\\di fluido in equilibrio.};
\end{tikzpicture}
\end{minipage}\hfill
\begin{minipage}[c]{0.60\textwidth}
\textbf{$F_0=p_0S$, verso il basso.} L'aria sopra il liquido esercita
questa forza sulla \textbf{faccia superiore} della colonna, distribuita
su tutta la sua area $S$.

\textbf{$F_P^{f}=m_fg$, verso il basso.} La Terra attrae tutto il fluido
della colonna: il peso agisce sull'\textbf{intero volume}. Nel disegno
è rappresentato da una sola freccia, applicata al baricentro.

\textbf{$R$, verso l'alto.} Il fluido sottostante esercita questa
forza di pressione sulla \textbf{faccia inferiore} della colonna,
distribuita su tutta la sua area $S$.

Il fluido circostante preme anche sulle facce laterali;
queste forze sono orizzontali e si bilanciano.
La colonna è ferma: la forza verso l'alto bilancia le due verso il basso.
\[
 R-F_0-F_P^{f}=0
 \quad\Longrightarrow\quad
 \boxed{R=F_0+F_P^{f}}.
\]
\end{minipage}

\vspace{6pt}
{\large\bfseries 2. Dalla forza sulla base alla pressione}

La pressione $p$ alla base si calcola dividendo $R$, la forza esercitata
\textbf{su quella faccia}, per la sua area $S$.
Il volume della colonna è $V=Sh$ e la sua massa è $m_f=d_fV$.
\begin{align*}
 p &= \frac{R}{S}
 && \text{definizione di pressione sulla base}\\[3pt]
   &= \frac{F_0+F_P^{f}}{S}
 && \text{per l'equilibrio della colonna}\\[3pt]
   &= \frac{F_0}{S}+\frac{F_P^{f}}{S}
 && \text{separiamo i due contributi}\\[3pt]
   &= p_0+\frac{m_fg}{S}
 && F_0=p_0S,\quad F_P^{f}=m_fg\\[3pt]
   &= p_0+\frac{d_fVg}{S}
 && \text{sostituiamo la massa }m_f=d_fV\\[3pt]
   &= p_0+\frac{d_fShg}{S}
 && \text{sostituiamo il volume }V=Sh\\[3pt]
   &= p_0+d_fgh
 && \text{semplifichiamo l'area }S.
\end{align*}

\vspace{5pt}
\begin{minipage}[c]{0.57\textwidth}
{\large\bfseries 3. Come varia la pressione?}
\[
  \boxed{\;p=p_0+d_fgh\;}
\]
Alla profondità $h=0$ la pressione è $p_0$.
Scendendo, aumenta linearmente: il grafico è una retta
con intercetta $p_0$ e coefficiente angolare $d_fg$.

È l'\textbf{aumento di pressione} $p-p_0=d_fgh$, rispetto alla faccia
superiore, a essere legato al peso della colonna.
\end{minipage}\hfill
\begin{minipage}[c]{0.39\textwidth}
\centering
\begin{tikzpicture}[x=1cm,y=1cm,>=Latex,font=\small]
  \draw[->] (-0.25,0) -- (4.15,0) node[below] {$h$};
  \draw[->] (0,-0.2) -- (0,3.15) node[left] {$p$};
  \node[below left] at (0,0) {$0$};
  \draw[very thick] (0,0.85) -- (3.65,2.8575);
  \fill (0,0.85) circle (1.6pt);
  \node[left] at (0,0.85) {$p_0$};
  \draw[dashed] (0.65,1.2075) -- (2.65,1.2075)
    node[midway,below] {$\Delta h$};
  \draw[dashed] (2.65,1.2075) -- (2.65,2.3075)
    node[midway,right] {$\Delta p$};
  \node[below] at (2,-0.5) {$\displaystyle\frac{\Delta p}{\Delta h}=d_fg$};
\end{tikzpicture}
\end{minipage}
\end{document}
